% Generated draft: requires mathematical and source review.
\documentclass[11pt]{article}
\usepackage[margin=25mm]{geometry}
\usepackage{fontspec}
\setmainfont{Latin Modern Roman}
\usepackage{amsmath,amssymb,mathtools,unicode-math}
\setmathfont{Latin Modern Math}
\usepackage{enumitem,needspace,xurl,hyperref}
\setlength{\parindent}{0pt}
\setlength{\parskip}{5pt}
\setlength{\emergencystretch}{3em}
\newcommand{\N}{\mathbb{N}}
\newcommand{\Z}{\mathbb{Z}}
\newcommand{\Q}{\mathbb{Q}}
\newcommand{\R}{\mathbb{R}}
\newcommand{\C}{\mathbb{C}}
\newcommand{\F}{\mathbb{F}}
\newcommand{\A}{\mathbb{A}}
\newcommand{\PP}{\mathbb{P}}
\newcommand{\E}{\mathbb{E}}
\newcommand{\eps}{\varepsilon}
\newcommand{\dd}{\,\mathrm{d}}
\begin{document}
\section*{UnsolvedMath ID: 30001275}
\textbf{Source problem code:} catalogue entry
\subsection{Definitions and mathematical statement}
Let \(E/\Q\) be an elliptic curve, let \(p\) be a prime at which \(E\) has good ordinary or multiplicative reduction, and let \(k\) be an imaginary quadratic field. Let \(k^{\mathrm{cyc}}/k\) and \(k^{\mathrm{ac}}/k\) denote its cyclotomic and anticyclotomic \(\Z_p\)-extensions, respectively, and put
\[
k_\infty=k^{\mathrm{cyc}}k^{\mathrm{ac}},\qquad G=\operatorname{Gal}(k_\infty/k),\qquad \Lambda(G)=\Z_p[[G]]=\varprojlim_U\Z_p[G/U],
\]
where \(U\) ranges over the open subgroups of \(G\).

For an explicit conventional interpretation of the dual Selmer group, let \(E[p^\infty]=\bigcup_{n\geq1}E[p^n]\). For each finite extension \(F/k\) contained in \(k_\infty\), define the usual \(p\)-primary Selmer group by
\[
\operatorname{Sel}_{p^\infty}(E/F)=\ker\!\left(H^1(F,E[p^\infty])\longrightarrow\prod_v H^1(F_v,E)[p^\infty]\right),
\]
where \(v\) ranges over all places of \(F\), the maps are induced by localization and the Kummer sequence, and \(H^1\) denotes continuous Galois cohomology. Set
\[
\operatorname{Sel}_{p^\infty}(E/k_\infty)=\varinjlim_F\operatorname{Sel}_{p^\infty}(E/F),\qquad
X(E/k_\infty)=\operatorname{Hom}_{\mathrm{cont}}\!\left(\operatorname{Sel}_{p^\infty}(E/k_\infty),\Q_p/\Z_p\right),
\]
using restriction maps in the direct limit and the discrete topology on the Selmer group. The dual has the induced \(\Lambda(G)\)-module structure. This is a stated conventional interpretation: the supplied excerpt does not specify its Selmer local conditions or whether any local conditions are omitted or modified.

Write \(f\) for the weight-two modular form associated with \(E\), as suggested by the notation in the source, and retain \(L_p(f,k_\infty)\) as the source's symbol for the intended two-variable \(p\)-adic \(L\)-function. Its construction, normalization, coefficient ring, and Euler-factor conventions are not supplied. For a finitely generated torsion \(\Lambda(G)\)-module \(M\), write \(\operatorname{char}_{\Lambda(G)}(M)\) for its characteristic ideal; for an element \(a\in\Lambda(G)\), write \((a)\) for the principal ideal it generates.

\textbf{Source conjectural target.} For every triple \((E,p,k)\) satisfying the stated hypotheses, is \(X(E/k_\infty)\) torsion over \(\Lambda(G)\), and does
\[
\bigl(L_p(f,k_\infty)\bigr)=\operatorname{char}_{\Lambda(G)}\!\bigl(X(E/k_\infty)\bigr)
\]
hold as an equality of ideals of \(\Lambda(G)\)? The characteristic-ideal assertion presupposes finite generation of \(X(E/k_\infty)\), and the left-hand side presupposes a specified integral realization of \(L_p(f,k_\infty)\) in that ring. Neither prerequisite is explained in the supplied excerpt. Consequently, this records the intended conjectural assertion, with the Selmer and analytic conventions still requiring confirmation. No restriction to odd \(p\), splitting hypothesis on \(p\) in \(k\), distinction between split and nonsplit multiplicative reduction, or additional conductor hypothesis is supplied; none is imposed here.
\subsection{Short English statement}
For an elliptic curve over \(\Q\) with good ordinary or multiplicative reduction at \(p\), the source asks whether its dual Selmer group over the combined cyclotomic and anticyclotomic extension of an imaginary quadratic field is torsion, with characteristic ideal generated by the intended two-variable \(p\)-adic \(L\)-function. The exact Selmer and \(L\)-function conventions remain unresolved. The supplied record labels the problem partially solved; this draft does not independently verify that assessment.
\subsection{Sources}
\begin{itemize}
\item Supplied local UnsolvedMath registry record, permanent ID \(30001275\), source problem code \texttt{catalogue entry}, title ``Cyclotomic--Anticyclotomic Iwasawa Main Conjecture,'' in the Oberwolfach Reports Open Problems collection. The record gives the proposed year as 2009 and names no individual proposer.
\item Supplied primary report citation: \emph{Algebraische Zahlentheorie} (2009), Oberwolfach Reports, pp.\ 1667--1730, DOI \texttt{10.4171/owr/2009/30}. \url{https://doi.org/10.4171/owr/2009/30}. The registry labels its statement \texttt{corrected\_verified} and reports checking Conjecture 1.5 in the official report, restoring the elliptic-curve, reduction, imaginary-quadratic-field, and two-extension hypotheses. This verification claim is supplied provenance and has not been independently checked here.
\item Supplied literature source labelled ``Rubin main conjecture'': \url{https://eudml.org/doc/143852}.
\item Supplied literature source labelled ``CM anticyclotomic Iwasawa theory'': \url{https://arxiv.org/abs/math/0302319}.
\item Supplied literature source labelled ``original OWR report'': \url{https://doi.org/10.4171/owr/2009/30}.
\item Literature assessment dated exactly \texttt{2026-08-21}: the supplied curation labels the problem partially solved and states that Rubin and Yager prove important CM cases of the Iwasawa main conjecture, while the general dihedral and two-variable statement in the extracted context remains only partially proved. This is supplied dated curation, not an independently verified literature assessment. The background explicitly calls for independent verification of time-sensitive status and cited claims.
\end{itemize}
\end{document}
